% ===================================================================
% SKILL: STANDART ZEBRA DESENLİ TABLO
% ===================================================================
% TANIM:
% Bu şablon, verileri sunmak için okunaklı, şık ve profesyonel bir 
% tablo yapısı oluşturur.
% 
% ÖZELLİKLER:
% 1. "booktabs" paketi (\toprule, \midrule, \bottomrule) ile modern çizgi yapısı.
% 2. "rowcolors" ile zebra deseni (okumayı kolaylaştırır).
% 3. "resizebox" ile tabloyu slayt genişliğine (veya istediğiniz orana) 
%    otomatik sığdırma.
% ===================================================================

\begin{frame}{Tablo Başlığı}
    % Tabloyu ortalar
    \centering 
    
    % Tablo Başlığı (Opsiyonel)
    \textbf{Tablo: Veri Dağılımı Örneği}
    \vspace{0.2em}
    
    % \resizebox{GENİŞLİK}{!}{İÇERİK}
    % Tablo çok büyükse 0.9\textwidth ile slayta sığdırır.
    \resizebox{0.9\textwidth}{!}{
        
        % Zebra Deseni: 2. satırdan başla, gri!10 ve beyaz renkleri değiştir.
        \rowcolors{2}{gray!10}{white}
        
        \begin{tabular}{c c c} % Sütun hizalamaları (c: center, l: left, r: right)
            \toprule
            \textbf{Başlık 1} & \textbf{Başlık 2} & \textbf{Başlık 3} \\
            \midrule
            Veri A-1          & Veri A-2          & Veri A-3          \\
            Veri B-1          & Veri B-2          & Veri B-3          \\
            Veri C-1          & Veri C-2          & Veri C-3          \\
            Veri D-1          & Veri D-2          & Veri D-3          \\
            \midrule
            \textbf{Toplam}   & \textbf{Sonuç}    & \textbf{\%100}    \\
            \bottomrule
        \end{tabular}
    }
\end{frame}
